\documentclass[12pt,a4paper]{article}

\usepackage[utf8]{inputenc}
\usepackage[T1]{fontenc}
\usepackage[margin=2.5cm]{geometry}
\usepackage{hyperref}
\usepackage{microtype}
\usepackage{parskip}
\usepackage{lmodern}
\usepackage{amsmath}

\hypersetup{colorlinks=true, urlcolor=blue!60!black}

\pagestyle{empty}

\begin{document}

\begin{flushright}
Kundai Farai Sachikonye \\
Auf der Lichtung 19, 82178 Puchheim, Germany \\
\href{https://github.com/fullscreen-triangle}{github.com/fullscreen-triangle} \\
\texttt{kundai.sachikonye@bitspark.com} \\
(+49) 1626386344 \\[6pt]
18 May 2026
\end{flushright}

\vspace{1em}

Prof.\ Dr.\ Detlef Schoder \\
Cologne Institute for Information Systems (CIIS) \\
Universität zu Köln \\
\texttt{schoder@wim.uni-koeln.de} \\
Ref.\ Wiss2604-xx

\vspace{1.5em}

\textbf{Re: Doctoral Researcher --- AI-Driven Protein Design (Ref.\ Wiss2604-xx)}

\vspace{1em}

Dear Prof.\ Schoder,

I am applying for the doctoral position in AI-driven protein design at CIIS
(Ref.\ Wiss2604-xx).
The position's core requirement --- integrating physics-informed scoring
functions into ML protein design pipelines and building hybrid scoring
frameworks that combine AI-derived and physics-based features --- is the
specific combination I have developed independently across several deployed
frameworks. The physics-based scoring is not a planned addition to my
toolkit; it is the foundation from which I approach molecular recognition problems.

\textbf{Physics-informed scoring for protein--protein interactions.}
The most important result I bring to the project is a zero-parameter method
for computing binding affinity from molecular structure.
The \href{https://github.com/fullscreen-triangle/crown-prince}{Crown Prince}
framework derives dissociation constants as:
\[
K_d = \exp\!\left(-\Delta M \cdot \ln b\right)
\]
where $\Delta M$ is the partition geometry distance between receptor and ligand
bounded oscillatory systems, and $b$ is the base of the ternary address encoding.
The Hill equation emerges as the exact limiting case of this expression when
receptor cooperativity collapses to a single binding mode.
Validated against 39 NIST binding compounds at zero free parameters.
This is a physics-informed scoring function in the precise sense the position
describes: a score derived from molecular geometry rather than from training data,
directly integrable into an ML protein design pipeline as a reward signal or
as a physics-based feature layer.

For protein--protein interaction prediction specifically, the
\href{https://syndrome-seven.vercel.app/tools/mhc-binding}{Syndrome}
framework applies spectral DFT to three physicochemical channels
(hydropathy, volume, charge) to compute a 36-dimensional cosine similarity
$\rho$ between any two molecular sequences.
Binding is predicted when $\rho \geq \rho^* = 0.72$; the binding ball
radius in spectral space is $r^* = \sqrt{2(1-\rho^*)} = 0.748$.
Validated: MHC-I/II binding prediction AUC\,$=0.81$ vs.\ NetMHCpan\,$=0.68$
on 2,847 IEDB peptides ($r=0.73$, $p<10^{-12}$).
The spectral complementarity $\rho$ is a physically grounded similarity
measure --- not a learned distance function --- that can be used directly
as a physics-based feature in a hybrid scoring framework, alongside
GNN-derived or AlphaFold-derived structural features.
For PPI lead optimization, the escape predictor tool screens all single
amino acid substitutions of a candidate binder for the variant that maximises
$\rho$ against the target while minimising off-target complementarity ---
a systematic computational lead optimization pass in $\mathcal{O}(cL\log L)$.

\textbf{ML pipeline architecture for protein design.}
The \href{https://honjo-masamune.vercel.app/tools}{Borgia} framework
generates molecular features --- electronic transitions, absorption/emission
spectra, vibrational modes, and partition coordinates --- from molecular
structure with zero empirical parameters, via the Triple Equivalence Theorem
(oscillatory dynamics $=$ categorical state enumeration $=$ partition operations).
These features serve directly as physics-informed input representations to
a protein design ML pipeline: instead of learning atomic embedding vectors
from a training set, the model receives features derived from molecular physics,
which generalise to out-of-distribution protein sequences by construction.
The \href{https://github.com/fullscreen-triangle/homo-veloce}{homo-veloce}
surrogate model system demonstrates the integration pattern in practice:
physics-based solvers (Gaussian processes, GNNs, ODE integrators) distilled
into lightweight 1D CNNs, LSTMs, and Transformers for real-time inference.
This is the same pattern as training an ML protein design model on
physics-informed features: the physics ensures structural validity;
the ML provides the generative and optimisation capability.
The \href{https://github.com/fullscreen-triangle/purpose}{Backward Training Principle}
formalises the training strategy: a binder model trained at the
\emph{pure-intent regime} --- the minimum binding energy configuration under
maximum on-target specificity, with no solubility or stability constraints ---
is $\epsilon$-competent at every constrained design sub-task (solubility, stability,
off-target avoidance) by feasibility inclusion, not generalisation.
A Forward Training Collapse Theorem establishes a positive VC failure gap
for the conventional approach of training on filtered feasible designs;
a Sample Complexity Theorem gives $(N{+}1)\times$ efficiency gains for a
cascade of $N$ design constraints. This directly addresses the sample
efficiency bottleneck in protein design benchmarks.

\textbf{Structural biology and protein dynamics.}
The \href{https://shakespear-nine.vercel.app/instrument}{Shakespear}
platform provides protein conformation, trajectory, catalysis, and dynamics
analysis via Kuramoto oscillator synchronization in S-entropy coordinate space.
Protein conformational states are represented as bounded oscillatory systems
with S-entropy coordinates $(S_k, S_t, S_e) \in [0,1]^3$; trajectory
completion reconstructs the full folding pathway from sparse observations
using a Banach contraction operator, without forward simulation.
For protein design benchmarks, this provides a lightweight computational
model for evaluating whether a designed binder adopts the predicted
fold --- before running full-scale MD or AlphaFold prediction.

\textbf{Agentic AI and ML pipeline engineering.}
The position lists willingness to engage with agentic AI approaches.
I have built production agentic systems: the
\href{https://github.com/fullscreen-triangle/four-sided-triangle}{Four-Sided-Triangle}
8-stage metacognitive RAG pipeline with Rust backend (10--50$\times$ speedup),
the \href{https://github.com/fullscreen-triangle/bloodhound}{Bloodhound}
distributed categorical VM with Triangle language (navigate/slice/complete),
and the \href{https://github.com/fullscreen-triangle/zangalewa}{Zangalewa}
Minimum Sufficient Interceptor for $O(\log_k N)$ specialist cascade routing.
For protein design, agentic architectures are relevant at the experimental design
layer: an agent that iterates between sequence proposal (generative ML),
physics scoring (Crown Prince $K_d$, spectral $\rho$), structural validation
(Shakespear trajectory completion), and experimental feedback is the pipeline
the position describes.

\textbf{Background.}
I hold an MSc in Computational Biology (Universität Tübingen) and conducted
doctoral research in Computational Lipidomics at TUM. I am legally resident
in Germany with full work authorisation. English is native; German is
conversational.

\vspace{1.5em}
Yours sincerely,

\vspace{2.5em}
\textbf{Kundai Farai Sachikonye}

\end{document}
